\subsection{Variable Neighborhood Search Mechanism}
\label{subsec:vns}
\label{subsec:vns_diversification}

If \(\mathrm{SampleSizeLocalSearch}\) does not improve the incumbent,
Algorithm~\ref{alg:framework} calls \(\mathrm{VNS}\) on the incumbent
chamber assignment and sample vector (Step~\ref{step:fw-vns}).
Following the classical VNS principle of changing neighborhood
intensity after local stagnation \citep{LiaoCheng2007}, and unlike
conventional VNS approaches that perturb sequencing or assignment
decisions, this mechanism perturbs sample-size decisions and therefore
changes the workload structure itself.

We present the main steps of the proposed VNS procedure in
Algorithm~\ref{alg:vns}.
It takes the incumbent solution as the initial solution.
The shake intensity is controlled by two destroy ratios,
\[
R=[0.20,0.30].
\]
The smaller ratio is tried first: it changes enough sample counts to
leave the current local basin, while still keeping much of the
incumbent allocation.
If that phase finds no improvement, the larger ratio destroys a
greater share of the sample vector.
For each ratio \(r\in R\) and for \(K_{\max}\) iterations, the
algorithm visits the two neighborhoods sequentially
(Steps~\ref{VNS:step1}--\ref{VNS:step7}).
In each neighborhood, \(\mathrm{Shake}\) selects a proportion \(r\) of
products and changes their sample sizes.
The first neighborhood (\(l=1\)) decreases the number of samples of
the selected products to \(s_{\min}\), while the second neighborhood
(\(l=2\)) increases the number of samples to \(v_{\max}\)
(Step~\ref{VNS:step2}).
The perturbed solution is then improved using the sample-size local
search in Algorithm~\ref{alg:sample_ls}.
The VNS follows a first-improvement policy: as soon as a locally
improved candidate improves the incumbent under the lexicographic
objective, it is returned to the main framework.
If no improvement is found within the shake budget of a ratio phase,
the method proceeds to the next destroy ratio.
If both phases fail, the incumbent is returned.

\begin{algorithm}[H]
\caption{VNS}
\label{alg:vns}
\label{alg:vns_fixed_destroy}
\label{alg:vns_short}
\begin{algorithmic}[1]
\State \textbf{Input:} chamber assignment \(\mathbf{x}\), incumbent solution \((\mathbf{s}_{inc},\Pi_{inc},T^{tot}_{inc})\)
\State \textbf{Output:} \((\mathbf{s},\Pi,T^{tot})\)
\State \(R\leftarrow[0.20,0.30]\)
\For{each \(r\in R\)}
    \For{\(k\leftarrow 1\) to \(K_{\max}\)}
        \For{\(l\leftarrow 1\) to \(2\)} \label{VNS:step1}
            \State \(\mathbf{s}_{cand}\leftarrow\mathrm{Shake}(\mathbf{s}_{inc},l,r)\) \label{VNS:step2}
            \State \((\mathbf{s}_{cand},\Pi_{cand},T^{tot}_{cand})\leftarrow\mathrm{SampleSizeLocalSearch}(\mathbf{x},\mathbf{s}_{cand})\)
            \If{\((T^{tot}_{inc},\mathbf{1}^{\mathsf{T}}\mathbf{s}_{inc})\succ_{\mathrm{lex}}(T^{tot}_{cand},\mathbf{1}^{\mathsf{T}}\mathbf{s}_{cand})\)}
                \State \Return \((\mathbf{s}_{cand},\Pi_{cand},T^{tot}_{cand})\)
            \EndIf
        \EndFor \label{VNS:step7}
    \EndFor
\EndFor
\State \Return \((\mathbf{s}_{inc},\Pi_{inc},T^{tot}_{inc})\)
\end{algorithmic}
\end{algorithm}

Together, the chamber-assignment, sample-size optimization, scheduling, and VNS components form an iterative framework in which workload generation, resource configuration, and schedule feasibility are updated jointly. The next section evaluates the contribution of these components through computational experiments.
